\section{Математическое ожидание.}
\subsection{Теоретическая справка.}
\begin{theoremdefinition}
    Следующие три определения математического ожидания эквивалентны:
    \begin{itemize}
        \item Пусть на вероятностном пространстве $(\Omega, \mathcal{F}, \P)$ задана случайная величина $\xi$. Математическим ожиданием $\xi$ называется интеграл Лебега $\Ex \xi = \int_\Omega \xi d\P$.
        \item Пусть на вероятностном пространстве $(\Omega, \mathcal{F}, \P)$ задана случайная величина $\xi$. Математическим ожиданием $\xi$ называется $\Ex \xi = \int_\R xd\P_\xi$.
        \item Пусть на вероятностном пространстве $(\Omega, \mathcal{F}, \P)$ задана случайная величина $\xi$ с распределением $F$. Математическим ожиданием $\xi$ называется интеграл Лебега-Стилтьеса вида $\Ex \xi = \int_\R xdF_{\xi}$.
    \end{itemize}
\end{theoremdefinition}
\begin{note}
    Поскольку математическое ожидание -- по сути просто интеграл Лебега, все теоремы, верные для него, верны и для математического ожидания.
\end{note}
\begin{proposition}
    Если $\xi$ и $\eta$~---~независимые случайные величины, то $\Ex(\xi \eta) \Ex \xi \Ex \eta$.
\end{proposition}
\begin{definition}
    Если существует математическое ожидание, то дисперсией называется $\D \xi = \Ex (\xi - \Ex \xi)^2$.
\end{definition}
\begin{definition}
    Ковариацией пары случайных величин называется $\cov(\xi, \eta) = \Ex(\xi - \Ex \xi)(\eta - \Ex \eta)$.
\end{definition}
\begin{definition}
    Моментом $k$-го порядка от случайной величины $\xi$ называется $\Ex \xi^k$.
    Абсолютным $k$-м моментом называется $\Ex |\xi|^k$, центральным -- $\Ex(\xi - \Ex \xi)^k$.
\end{definition}